\subsection{WEIGHTS OF A FINITE DIMENSIONAL SIMPLE g-MODULE}

PROPOSITION 1. Let V be a finite dimensional g-module.

\begin{enumerate}
\def\labelenumi{(\roman{enumi})}
\item
  All the weights of V belong to P.
\item
  \(V = \bigoplus_{\mu \in P} V^{\mu}.\)
\item
  For all \(\mu \in \mathfrak{h}^*\) , \(\mathrm{V}^{\mu}\) is the set of \(x \in \mathrm{V}\) such that \(h.x = \mu(h)x\) for all \(h \in \mathfrak{h}\) .
\end{enumerate}

For all \(\alpha \in \mathrm{B}\) , there exists a homomorphism from \(\mathfrak{sl}(2,k)\) to \(\mathfrak{g}\) that takes \(H\) to \(H_{\alpha}\) . Thus, by §1, no. 2, Cor. of Prop. 2, \((H_{\alpha})_{\mathrm{V}}\) is diagonalizable and its eigenvalues are integers. Hence, the set of \((H_{\alpha})_{\mathrm{V}}\) , for \(\alpha \in \mathrm{B}\) , is diagonalizable (Algebra, Chap. VII, §5, no. 6, Prop. 13). Consequently, for all \(h \in \mathfrak{h}\) , \(h_{\mathrm{V}}\) is diagonalizable. By Chap. VII, §1, no. 3, Prop. 9, \(\mathrm{V} = \bigoplus_{\mu \in \mathfrak{h}^*} \mathrm{V}^{\mu}\) . On the other hand, if \(\mathrm{V}^{\mu} \neq 0\) , the preceding shows that \(\mu(H_{\alpha}) \in \mathbf{Z}\) for all \(\alpha \in \mathrm{B}\) , so \(\mu \in \mathrm{P}\) . This proves (i) and (ii). We see in the same way that \(\mathfrak{h}_{\mathrm{V}}\) is diagonalizable, hence (iii).

COROLLARY. Let \(\rho\) be a finite dimensional representation of \(\mathfrak{g}\) and \(\Phi\) the bilinear form associated to \(\rho\) .

\begin{enumerate}
\def\labelenumi{(\roman{enumi})}
\item
  If \(x, y \in \mathfrak{h}_{\mathbf{Q}}\) , then \(\Phi(x, y) \in \mathbf{Q}\) and \(\Phi(x, x) \in \mathbf{Q}_{+}\) .
\item
  If \(\rho\) is injective, the restriction of \(\Phi\) to \(\mathfrak{h}\) is non-degenerate.
\end{enumerate}

Assertion (i) follows from Prop. 1 since the elements of P have rational values on \(\mathfrak{h}_{\mathbf{Q}}\) . If \(\rho\) is injective, \(\Phi\) is non-degenerate (Chap. I, §6, no. 1, Prop. 1), so the restriction of \(\Phi\) to \(\mathfrak{h}\) is non-degenerate (Chap. VII, §1, no. 3, Prop. 10 (iii)).

Lemma 1. Let \(\mathrm{V}\) be a \(\mathfrak{g}\) -module and \(\rho\) the corresponding representation of \(\mathfrak{g}\) .

\begin{enumerate}
\def\labelenumi{(\roman{enumi})}
\tightlist
\item
  If \(a\) is a nilpotent element of \(\mathfrak{g}\) , and if \(\rho(a)\) is locally nilpotent,
\end{enumerate}

\[
\rho (e ^ {\mathrm{ad} a} b) = e ^ {\rho (a)} \rho (b) e ^ {- \rho (a)}
\]

for all \(b \in \mathfrak{g}\).

\begin{enumerate}
\def\labelenumi{(\roman{enumi})}
\setcounter{enumi}{1}
\tightlist
\item
  If \(\alpha \in \mathrm{R}\) and if the images under \(\rho\) of the elements of \(\mathfrak{g}^{\alpha}\) and \(\mathfrak{g}^{-\alpha}\) are locally nilpotent, the set of weights of \(V\) is stable under the reflection \(s_{\alpha}\) .
\end{enumerate}

With the assumptions in (i), we have \(\rho((\mathrm{ad}a)^n b) = (\mathrm{ad}\rho(a))^n\rho(b)\) for all \(n\geq 0\) , so \(\rho(e^{\mathrm{ad}a}b) = e^{\mathrm{ad}\rho(a)}\rho(b)\) . On the other hand,

\[
e ^ {\mathrm{ad} \rho (a)} \rho (b) = e ^ {\rho (a)} \rho (b) e ^ {- \rho (a)}
\]

is assertion (ii) of Chap. VII, §3, no. 1, Lemma 1.

We now adopt the assumptions in (ii). Let \(\theta_{\alpha} = e^{\mathrm{ad}X_{\alpha}}e^{\mathrm{ad}X_{-\alpha}}e^{\mathrm{ad}X_{\alpha}}\) . By (i), there exists \(S\in \mathbf{GL}(V)\) such that \(\rho (\theta_{\alpha}b) = S\rho (b)S^{-1}\) for all \(b\in \mathfrak{g}\) . Now \(\theta_{\alpha}|\mathfrak{h}\) is the transpose of \(s_{\alpha}\) (§2, no. 2, Lemma 1). Let \(\lambda\) be a weight of V. There exists a non-zero element \(x\) of V such that \(\rho (h)x = \lambda (h)x\) for all \(h\in \mathfrak{h}\) . Then

\[
\rho (h) \mathrm{S} ^ {- 1} x = \mathrm{S} ^ {- 1} \rho (^ {t} s _ {\alpha} h) x = \mathrm{S} ^ {- 1} \lambda (^ {t} s _ {\alpha} h) x = (s _ {\alpha} \lambda) (h) \mathrm{S} ^ {- 1} x
\]

for all \(h \in \mathfrak{h}\) . Consequently, \(s_{\alpha}\lambda\) is a weight of V.

PROPOSITION 2. Let \(\mathrm{V}\) be a finite dimensional \(\mathfrak{g}\) -module and \(s \in \operatorname{Aut}_0(\mathfrak{g})\) .

\begin{enumerate}
\def\labelenumi{(\roman{enumi})}
\item
  There exists \(\mathrm{S} \in \mathbf{GL}(\mathrm{V})\) such that \((s(x))_{\mathrm{V}} = \mathrm{S}x_{\mathrm{V}}\mathrm{S}^{-1}\) for all \(x \in \mathfrak{g}\) .
\item
  If \(s \in \mathrm{Aut}_e(\mathfrak{g})\) , \(S\) can be chosen to be an element of \(\mathbf{SL}(V)\) leaving stable all the \(\mathfrak{g}\) -submodules of \(V\) .
\end{enumerate}

Assertion (ii) follows from Lemma 1 (i). Now let \(s \in \mathrm{Aut}_0(\mathfrak{g})\) and denote by \(\rho\) the representation of \(\mathfrak{g}\) defined by V. By (ii), the representations \(\rho\) and \(\rho \circ s\) become equivalent after extension of scalars. They are therefore equivalent (Chap. I, §3, no. 8, Prop. 13), hence the existence of S.

Remark 1. Let S satisfy the condition in Prop. 2 (i), and let \(\mathfrak{h}' = s(\mathfrak{h})\) ; denote by \(s^*\) the isomorphism \(\lambda \mapsto \lambda \circ s^{-1}\) from \(\mathfrak{h}^*\) to \(\mathfrak{h}'^*\) . It is clear that

\[
\mathrm{S} (\mathrm{V} ^ {\lambda}) = \mathrm{V} ^ {s ^ {*} \lambda}.
\]

In particular:

COROLLARY 1. The isomorphism \(s^*\) takes the weights of V with respect to \(\mathfrak{h}\) to those of V with respect to \(\mathfrak{h}'\) ; corresponding weights have the same multiplicity.

COROLLARY 2. Let \(w \in \mathrm{W}\) . For all \(\lambda \in \mathfrak{h}^*\) , the vector subspaces \(\mathrm{V}^{\lambda}\) and \(\mathrm{V}^{w\lambda}\) have the same dimension. The set of weights of \(\mathrm{V}\) is stable under \(\mathrm{W}\) .

Indeed, \(w\) is of the form \(s^*\) with \(s \in \mathrm{Aut}_e(\mathfrak{g}, \mathfrak{h})\) (§2, no. 2, Cor. of Th. 2).

Remark 2. By Cor. 1 of Prop. 2 and §5, no. 3, Remark 2, it makes sense to speak of the weights of V with respect to the canonical Cartan subalgebra of g, and of their multiplicities.

Remark 3. Lemma 1 (i) and Prop. 2 remain valid, with the same proof, even if g is not assumed to be splittable.

